\chapter{論文末の問いと見通し}
\label{ch:questions}

原論文は，次の問題を今後の課題として挙げている．
ここでは 1984 年当時の問題提起として記録する．

\begin{enumerate}
  \item 連続な周辺分布に対する最適結合は，
    直積測度に関して常に特異になるのか．
  \item 経験測度は $1<p<\infty$ に対して，どの条件のもとで
    $\Wass_p$ に関して概収束するか．
  \item 無限次元線形空間，特に Hilbert 空間上の Gaussian 測度について，
    $\Wass_2$ を明示的に計算できるか．
\end{enumerate}

\begin{remark}{有限次元公式の意味}{gs-formula-significance}
  共分散行列が可換な場合だけでなく，一般の非可換な場合にも
  \[
    \tr\!\left[
      \left(
        \Sigma_0^{1/2}\Sigma_1\Sigma_0^{1/2}
      \right)^{1/2}
    \right]
  \]
  が現れることが本論文の要点である．
  単純な差 $\norm{\Sigma_0^{1/2}-\Sigma_1^{1/2}}_{\mathrm{F}}$ は，
  共分散行列が可換な場合にしか正しい値を与えない．
\end{remark}

\begin{remark}{参考文献}{gs-reference}
  C. R. Givens and R. M. Shortt,
  \emph{A class of Wasserstein metrics for probability distributions},
  Michigan Math. J. \textbf{31} (1984), no. 2, 231--240.
  DOI: 10.1307/mmj/1029003026.
\end{remark}
